% Appendix B: Setup and Run Instructions.
% Sources: README.md (setup, data, tunnel, checks), cli.py at HEAD (flags), requirements-observability.txt
% (Phoenix), docs/final/B_status.md and FACTS.md section 17 (evaluation commands), FACTS.md section 16
% (versions, pip list / node / ollama on 2026-11-12), runs/eval_seed7_latest.json (manifest hash).
% No key or token value is printed anywhere in this appendix.
\chapter{Setup and Run Instructions}
\label{app:run}

The commands below match the repository \texttt{README.md}. The repository is \url{https://github.com/Suraj-B12/StripCheck}. StripCheck needs no account, no API key and no downloaded language model; the optional local model and the optional tracing server are described in Section~\ref{sec:app-run-optional}.

\section{Tested environment}
\label{sec:app-run-env}

Table~\ref{tab:app-versions} lists the versions on the Windows~11 laptop that produced the reported evaluation. The README asks for Python 3.12 or later and Node.js 22.12 or later; the build was also exercised on Python 3.14 under macOS.

\begin{table}[H]
\centering
\caption{Software versions of the evaluation machine}
\label{tab:app-versions}
\small
\begin{tabularx}{\linewidth}{@{}l >{\raggedright\arraybackslash}X@{}}
\toprule
\textbf{Component} & \textbf{Version} \\
\midrule
Python and API & Python 3.12.7; FastAPI 0.162.0; Uvicorn 0.54.0; Pydantic 2.14.1 \\
Pack reading & RapidOCR 3.10.0 (PP-OCRv6-small); ONNX Runtime 1.31.1; zxing-cpp 3.1.1; OpenCV (headless) 4.14.0.94; Pillow 13.0.0 \\
Matching and tests & RapidFuzz 3.14.6; pytest 9.1.1 \\
Web client & Node.js 24.11.1; npm 11.6.2; React 19.3.1; Vite 8.3.5; TypeScript 6.0.3; Vitest 5.0.5 \\
Optional planner model & Ollama 0.17.7 with \texttt{qwen2.5:3b} (1.9~GB) \\
Optional tracing & OpenTelemetry SDK and OTLP/HTTP exporter 1.45.1; Arize Phoenix 21.0.0 in a separate environment \\
\bottomrule
\end{tabularx}
\end{table}

\section{First setup on a fresh clone}
\label{sec:app-run-setup}

On Windows PowerShell:
\begin{lstlisting}[language={},numbers=none]
git clone https://github.com/Suraj-B12/StripCheck.git
cd StripCheck
py -3.12 -m venv .venv
.\.venv\Scripts\python.exe -m pip install -r requirements.txt -r requirements-dev.txt
npm.cmd --prefix web ci
npm.cmd --prefix web run build
.\.venv\Scripts\python.exe cli.py sync --full     # fetch every published month (several minutes)
.\.venv\Scripts\python.exe cli.py serve           # website and API on http://127.0.0.1:8000
\end{lstlisting}

On macOS or Linux:
\begin{lstlisting}[language={},numbers=none]
python3 -m venv .venv
source .venv/bin/activate
python -m pip install -r requirements.txt -r requirements-dev.txt
npm --prefix web ci
npm --prefix web run build
python cli.py sync --full
python cli.py serve
\end{lstlisting}

The pinned Python versions recorded on the macOS build machine can be applied with
\begin{lstlisting}[language={},numbers=none]
pip install -r requirements.txt -r requirements-dev.txt -c requirements.lock.txt
\end{lstlisting}
For later Windows launches, \texttt{./start-stripcheck.ps1} starts the server (add \texttt{-Lan} for token-protected phone pairing). If port 8000 is in use, run \texttt{python cli.py serve --port 8001}. API documentation is served at \texttt{/api/docs}.

\section{Alert data}
\label{sec:app-run-data}

\begin{lstlisting}[language={},numbers=none]
python cli.py sync            # refresh the newest three published months
python cli.py sync --full     # all published months (requests at least 1.5 s apart)
python cli.py sync --offline  # rebuild the index from stored snapshots, no network
python cli.py stats           # statistics of the local database
python cli.py doctor          # installed readers, local database, Ollama and Phoenix status
\end{lstlisting}

Each sync stores the exact source bytes with their SHA-256 checksum and link under \texttt{data/raw/portal/}. Repeating an unchanged import adds no duplicate rows, and rejected rows stay visible in the quarantine view. Photos, snapshots, the SQLite database and logs stay in the ignored \texttt{data/} and \texttt{runs/} directories; \texttt{STRIPCHECK\_DATA\_DIR}, \texttt{STRIPCHECK\_DB\_PATH} and \texttt{STRIPCHECK\_RUNS\_DIR} move them elsewhere.

Table~\ref{tab:ingest} gives the validation outcomes of the ingest described in
Section~\ref{sec:impl-data-preparation-and-indexing}, from a read-only query of the local database.
The quarantine reasons are listed as exclusive groups, because a row can carry more than one reason.

\begin{table}[htbp]
\centering
\caption[Validation outcomes in the local index]{Validation outcomes in the local index (read-only query of \texttt{data/stripcheck.db})}
\label{tab:ingest}
\footnotesize
\begin{tabularx}{\linewidth}{@{}>{\raggedright\arraybackslash}p{3.3cm} >{\raggedright\arraybackslash}X r@{}}
\toprule
\textbf{Outcome} & \textbf{Reason or flag} & \textbf{Rows} \\ \midrule
Quarantined (75 rows) & batch placeholder only (empty, too short, ``not mentioned'') & 47 \\
 & manufacturer missing on a non-spurious entry only & 14 \\
 & both batch placeholder and manufacturer missing & 8 \\
 & batch cell holds a date (shifted column) & 5 \\
 & several batches that cannot be split reliably & 1 \\ \midrule
Indexed with a flag & batch key under 4 characters / non-drug item & 144 / 98 \\
 & date missing / text encoding repaired & 97 / 29 \\
 & manufacture after list month / after expiry & 16 / 14 \\ \midrule
Active alerts & rows used by checks and by the evaluation & \ResAlertCount{} \\
\bottomrule
\multicolumn{3}{@{}p{\linewidth}@{}}{\footnotesize An indexed row can carry several flags, so the flag counts overlap.}
\end{tabularx}
\end{table}

\section{Checking a pack}
\label{sec:app-run-check}

The web client at \url{http://127.0.0.1:8000} accepts one or two photos or typed details. The same checks run from the command line:
\begin{lstlisting}[language={},numbers=none]
python cli.py check front.jpg back.jpg              # one or two photos, rule planner
python cli.py check front.jpg --planner hybrid      # local LLM at the branch points
python cli.py check --text LM2408B --product "..." --maker "..."
\end{lstlisting}

\section{Configuration, local model and tracing}
\label{sec:app-run-optional}

Settings are environment variables, listed in \texttt{.env.example}. A local copy named \texttt{.env} is read at start-up and must never be committed. No setting holds a third-party key. The pairing token for phone access (\texttt{STRIPCHECK\_LAN\_TOKEN}) is generated per session by the launch scripts and is not written into any file in the repository.

The hybrid planner needs a local Ollama server with the model pulled. Without it, every branch point falls back to the rules and the check completes.
\begin{lstlisting}[language={},numbers=none]
ollama pull qwen2.5:3b
$env:STRIPCHECK_PLANNER = "hybrid"      # default is "rules"
python cli.py doctor                    # reports whether Ollama is up and which models it has
\end{lstlisting}
A cold model load can take longer than the per-call timeout, in which case the first decision falls back to the rules; the walkthrough script therefore preloads the model before its first check.

Tracing to Arize Phoenix is off by default; the JSONL step log is always written.
\begin{lstlisting}[language={},numbers=none]
.\.venv\Scripts\python.exe -m pip install -r requirements-observability.txt
python -m venv .venv-phoenix
.\.venv-phoenix\Scripts\python.exe -m pip install arize-phoenix
.\.venv-phoenix\Scripts\phoenix.exe serve      # UI on http://127.0.0.1:6006
$env:STRIPCHECK_TRACING = "1"
\end{lstlisting}

\begin{sloppypar}The trace walkthrough of Appendix~\ref{app:traces} is reproduced by \texttt{docs/\allowbreak trace/\allowbreak run\_walkthrough.py}, which uses a temporary data folder and one invented alert record, so the real database is not touched.\end{sloppypar}

\section{Live phone demonstration}
\label{sec:app-run-tunnel}

\begin{lstlisting}[language={},numbers=none]
powershell -ExecutionPolicy Bypass -File .\scripts\start-demo-tunnel.ps1
\end{lstlisting}
The script starts the server on localhost, opens a temporary Cloudflare quick tunnel and prints a phone link that carries a pairing token generated for that session, together with a QR code. Requests without the token receive HTTP 401. Ctrl-C stops both processes. The script needs \texttt{cloudflared} and no account. The link should not be shared beyond the people taking part in the demonstration.

\section{Tests}
\label{sec:app-run-tests}

\begin{lstlisting}[language={},numbers=none]
python -m pytest -q           # runs without Ollama, Phoenix or internet
python cli.py selftest        # offline end-to-end run of the synthetic scenarios
python cli.py eval --quick    # the 11 authored diagnostic scenarios
npm --prefix web run build
npm --prefix web test         # Vitest
\end{lstlisting}
Test counts at the commit described are given in Chapter~\ref{ch:req}.

\section{Reproducing the evaluation}
\label{sec:app-run-eval}

The following commands regenerate the evaluation set and every number in Chapters~\ref{ch:eval} and~\ref{ch:exp}. The hybrid arm needs Ollama with \texttt{qwen2.5:3b}.
\begin{lstlisting}[language={},numbers=none]
.venv\Scripts\python.exe cli.py props --eval --n 60 --seed 7
.venv\Scripts\python.exe cli.py eval --arms baseline,agent_rules --hitl oracle,none --seed 7
.venv\Scripts\python.exe cli.py eval --arms agent_hybrid --hitl oracle,none --seed 7
\end{lstlisting}
\begin{sloppypar}The commands write \texttt{runs/eval\_<timestamp>.json} and \texttt{.md}, the merged \texttt{runs/eval\_seed7\_latest.json}, \texttt{docs/results/eval\_seed7.md}, and the LaTeX files \texttt{paper/generated/results.tex} and \texttt{results\_tables.tex} that this report reads. The development set uses \texttt{--seed 1}. The packs are rendered from rows of the local index, so a later sync can change the set; the reported run used \res{ResAlertCount} indexed alerts and the evaluation manifest with hash \texttt{73de16c2934fa4a0}. The report itself is built with \texttt{cd paper/report; latexmk main.tex}, with output in \texttt{paper/report/build/}.\end{sloppypar}

\section{Engine modules}
\label{sec:app-run-modules}

Table~\ref{tab:modules} lists the modules of the engine in \texttt{core/}, as referred to in
Section~\ref{sec:impl-repository-structure}.

\begin{table}[htbp]
\centering
\caption{Engine modules in \texttt{core/}}
\label{tab:modules}
\footnotesize
\begin{tabularx}{\linewidth}{@{}>{\ttfamily}l >{\raggedright\arraybackslash}X@{}}
\toprule
\normalfont\textbf{Module} & \textbf{Purpose} \\ \midrule
investigator.py & Bounded loop, branch points, questions, answers, state, span export \\
verifier.py & Rules R1 to R10; downgrade to \CANTV{} \\
matching.py, noisy.py & Exact and variant search, channel score, grading; confusion prior, positions to confirm \\
textutils.py & Batch key, look-alike table, dates, product coverage, maker names, GS1 parsing \\
templates.py & All verdict and question text, banned-phrase filter, citations \\
llm.py & Optional Ollama calls with schema, timeout, fallback, circuit breaker \\
pack\_tools.py, crops.py & Image validation, barcode and text readers, field extraction, character crops \\
portal.py, ingest.py & Portal sync and snapshots; cleaning, quarantine, idempotent import \\
supervisor.py & Sync, ingest and cabinet re-check; confusion memory \\
db.py, config.py, tracing.py & SQLite schema and documents; bounded settings; JSONL and OpenTelemetry traces \\
golden.py, props.py, evaluate.py & Synthetic scenarios and cartons; evaluation arms, metrics, generated tables \\
\bottomrule
\end{tabularx}
\end{table}
